% ID: FIS-TER-DIL-001
% Titolo: Foro in un tassello di rame riscaldato
% Disciplina: Fisica
% Area: Termologia
% Argomento: Dilatazione lineare
% Sottoargomento: Dilatazione di solidi con fori
% Classe: Seconda liceo scientifico
% Difficolta: 3
% Tipo: quesito_teorico
% Risultato: il foro aumenta di diametro
% Tempo_stimato: 6 min
% Competenze: interpretazione fisica, modellizzazione, argomentazione
% Prerequisiti: dilatazione lineare, dilatazione termica dei solidi
% Tag: termologia, dilatazione_lineare, rame, foro, solidi
% Fonte: esempio_originale
% Autore: Riccardo Carli
% Licenza: CC-BY-SA-4.0

\begin{esercizio}
Un tassello di rame presenta un piccolo foro circolare. Il tassello viene
riscaldato in modo uniforme. Il diametro del foro aumenta, diminuisce oppure
rimane invariato? Motiva la risposta.
\end{esercizio}

\begin{soluzione}
Il diametro del foro aumenta. Il modo piu semplice per capirlo e immaginare che
il foro sia riempito da un disco dello stesso materiale: se il tassello e il
disco venissero riscaldati insieme, entrambi si dilaterebbero nello stesso modo.
Il disco diventerebbe piu grande, quindi anche lo spazio che lo conterrebbe deve
diventare piu grande.

La dilatazione uniforme non chiude il foro: tutte le lunghezze caratteristiche
del corpo, comprese quelle che descrivono cavita e distanze interne, vengono
moltiplicate per lo stesso fattore di dilatazione.
\end{soluzione}

